% ===== Chapter 2: auto-converted from Word =====

\section{Literature Review}

The relationship between insider ownership and shareholder returns sits at the intersection of three large but partially connected literature pieces: the agency cost theory, the empirical literature on ownership concentration and firm value, and the asset pricing literature on ownership sorted portfolio returns. This chapter integrates these strands with two objectives. The first is to show that, despite nearly five decades of empirical work since Jensen and Meckling (1976), the literature on the risk adjusted return implications of high insider ownership remains geographically concentrated in the United States, and methodologically heterogeneous. The literature on the relationship between insider ownership and stock returns with a long sample and globally oriented remains thin. The second is to position the present thesis precisely against this gap, building on the descriptive global evidence assembled by Wood (2025) and the portfolio sorting framework introduced by Fahlenbrach (2009) and refined by Lilienfeld-Toal and Ruenzi (2014). Throughout the paper, the thought that motivates this thesis is a long horizon, global, factor adjusted analysis of insider ownership sorted portfolios. 


\subsection{From Agency Costs to Stewardship}

Modern theories of ownership and firm performance lead back to Berle and Means (1932). They brought to the light that the separation of ownership from control in corporations creates a fundamental misalignment. The interests of professional managers and those of the firm's shareholders. This concern was formalised by Jensen and Meckling (1976), whose logical model treats managers as utility maximizing agents whose private interests diverge from those of outside shareholders. This generates costs that include perquisite consumption, empire building, and effort aversion. In their framework, increasing the equity stake held by management reduces these agency costs by aligning managerial cashflow rights with those of outside shareholders. This is called the convergence of interest hypothesis. The theoretical prediction is unambiguous: at the margin, higher inside ownership should be associated with better firm performance.

This unambiguous prediction was complicated by Fama and Jensen (1983), who argued that very high concentration of voting and cashflow rights in the hands of insiders can entrench management beyond the reach of the labour and capital markets. Once an insider's stake exceeds the threshold at which external monitoring can threaten removal, agency costs may rise with ownership rather than fall. This as entrenched managers extract private benefits of control at the expense of minority shareholders. The interplay between these two opposing powers, alignment and entrenchment, generates the non-monotonic theoretical relationship between ownership and value that has dominated empirical testing for nearly four decades.

Another theoretical perspective, stewardship theory, was developed by Donaldson and Davis (1991) as a partial counterweight to the agency framework. Where agency theory views managers as self-interested utility maximizers, requiring incentive alignment. Donaldson and Davis (1991) through stewardship theory describe managers as those with significant ownership stakes, as intrinsically motivated to act as long term stewards of the firm. The stewardship perspective is discussed extensively by Wood (2025), who argues that the behavioural patterns observed in his global sample of high ownership firms have lower earnings guidance, higher reinvestment and lower dividend payout. These are inconsistent with pure agency theoretic predictions and more correctly explained by long horizon stewardship motives.

For the purposes of this thesis, these theoretical elements matter because they generate competing predictions about the return implications of high insider ownership. Pure alignment predicts that high ownership firms outperform and entrenchment predicts that very high ownership firms underperform. Stewardship predicts long run outperformance driven by reinvestment and value creation rather than short term profit extraction. Distinguishing among these empirically requires risk adjusted return data over a long horizon and across institutional environments.


\subsection{Insider Ownership and Firm Value}

The empirical literature on insider ownership has focused on firm value, typically measured by Tobin's Q, rather than on stock returns. The seminal contribution by Morck, Shleifer, and Vishny (1988) used a methodology working with ranges on 371 Fortune 500 firms in 1980 and documented the common and non-monotonic pattern: alignment effects dominate from 0\% to 5\% ownership, entrenchment effects dominate between 5\% and 25\%, and alignment effects emerge above 25\% again. McConnell and Servaes (1990) refined this finding using a substantially larger sample of 1133 firms across two cross sections (1976 and 1986). They found a smoother inverted U relationship, with Q rising until roughly 40-50\% ownership before declining. Hermalin and Weisbach (1991), working with a smaller longitudinal sample, finding similar non-monotonic patterns but with weaker statistical significance.

The cross sectional value literature took a methodologically important turn with Demsetz and Villalonga (2001), who argued that the entire ownership value relationship is fundamentally endogenous: ownership structures are themselves chosen optimally given the firm's environment and performance. Treating ownership as exogenous in a regression of Q on ownership produces biased estimates. Once ownership is instrumented, they reported, the statistical relationship between ownership and Q largely disappears. This endogeneity critique remains the most serious methodological objection to the entire ownership value literature and has shaped subsequent empirical design. 

Anderson and Reeb (2003) shifted the emphasis from generic insider ownership to the more specific category of family ownership in S\&P 500 firms during 1992-1999. Reporting a robust positive association between family control and firm value. Villalonga and Amit (2006) split family ownership into ownership, control, and management dimensions. Showing that the value premium attaches specifically to firms in which the founder remains active as CEO or chairman. Where family control through dual class shares or pyramidal structures destroys value. The latter finding is consistent with the entrenchment hypothesis.

A more recent and contrarian contribution is provided by Fabisik, Fahlenbrach, Stulz, and Taillard (2021), who use a 27 year panel of 1852 US firms. They find that, contrary to the earlier consensus, firms with higher managerial ownership are on average worth less in terms of Tobin's Q. They attribute this to selection: firms with weaker outside investor demand retain higher insider stakes by default. The apparent "ownership premium" in earlier studies reflects unobserved firm characteristics rather than a causal effect of ownership on value. This result is important for the this paper because it sharpens the distinction between value based and return based evidence. A firm can simultaneously trade at a lower Q (consistent with Fabisik et al.) and deliver higher risk adjusted returns going forward (consistent with Wood, 2025). The low Q reflects investor underpricing rather than worsening fundamentals. Risk adjusted return tests are precisely the core that distinguishes these interpretations.

What unites the value literature, despite its disagreements between the papers, is its focus on US data, its use of Tobin's Q as the dependent variable, and its reliance on point in time cross sections. As Wood (2025) observes in his survey of more than one hundred prior works, no peer reviewed study has combined long horizons, global coverage, and risk adjusted return measurement. This paper addresses that under investigated piece of literature.


\subsection{Insider Ownership and Stock Returns}

The literature directly relevant to this thesis is the literature on ownership and stock price performance. It is much smaller than the value literature and only a handful of papers dominate it. Before reviewing those papers, it is important to clarify a distinction that the literature often leaves out, but that matters for this thesis. CEO ownership, founder ownership, family ownership, and aggregate insider ownership are not the same thing. They capture different governance phenomena and may carry different return implications.


\subsubsection{Separating Ownership}

CEO ownership is the equity stake held by the chief executive officer alone. This is the variable used by Lilienfeld-Toal and Ruenzi (2014). It captures one individual's incentive alignment. It says nothing about the rest of the board or the broader executive team. Founder ownership, used by Fahlenbrach (2009) and Adams, Almeida, and Ferreira (2009), conditions on the founder still being active in the firm as CEO or chairman. It combines ownership with a specific identity and tenure dimension. Family ownership, used by Anderson and Reeb (2003), Villalonga and Amit (2006), Sraer and Thesmar (2007), Andres (2008), and Lins, Volpin, and Wagner (2013), captures the aggregate stake of a founding family. It typically requires identifying the family as a related ownership block and may include whose incentives differ from the founder's. Aggregate insider ownership, used by Wood (2025) and adopted in this thesis, sums the stakes of all members of the executive team and the board. It does not condition on whether those individuals are founders, family members, or unrelated managers.

These four constructs overlap but are not interchangeable. A firm with a high CEO ownership stake may or may not have high aggregate insider ownership, depending on the size of stakes held by the rest of the board and executive team. A founder led firm in which the founder holds 30\% but other insiders hold negligible stakes will register as high CEO ownership and high aggregate insider ownership. A mature firm in which a non founder CEO holds 5\% but the broader board and executive team collectively hold an additional 20\% will register as low CEO ownership but high aggregate insider ownership. Family firms may have high family ownership but low CEO ownership when the chief executive is a professional manager rather than a family member.

The empirical implications of these distinctions are not small. The theoretical mechanisms work differently across the different ownership definitions. CEO ownership most directly tests the Jensen and Meckling (1976) alignment hypothesis at the level of the individual decision maker. Founder ownership additionally captures the long horizon motivation emphasised by stewardship theory (Donaldson and Davis, 1991). Family ownership introduces dynastic and reputational considerations that do not apply to non family insiders. Aggregate insider ownership captures the collective stake of the entire control group. This is most directly relevant to questions of board level monitoring, collective decision making, and the alignment of the firm's governance with outside shareholders.


\subsubsection{Founder CEO Ownership: Fahlenbrach (2009)}

Fahlenbrach (2009) is the most frequently cited methodological precedent in this literature. He studies 361 US founder CEO firms over 1992-2003. He forms calendar time portfolios of founder CEO firms versus a matched non founder benchmark. He estimates four factor (Carhart, 1997) regressions on monthly portfolio returns. He documents a value weighted alpha of roughly 4.4\% per year for founder CEO firms. The alpha is robust to controls for firm characteristics, CEO characteristics and industry affiliation. Fahlenbrach interprets the alpha not as evidence of mispricing but as a return to the different investment behaviour of founder managed firms. The higher capital expenditure, more research and development, and more focused acquisition strategies. The behavioural mechanisms he identifies are shown almost exactly in the descriptive findings of Wood (2025). This suggests that the broader owner population may share the behavioural drivers that Fahlenbrach attributed specifically to founders.

Fahlenbrach's contribution is narrower than this thesis. Firms in his treatment group share the founder active condition but vary widely in actual ownership stake. His sample is restricted to US large- and midcap firms. This thesis differs in that it uses ownership in a broader sense, applying it to a global universe, and including firms in which the high ownership block is held by non founder insiders. This is a population that Fahlenbrach's methodology does not address.


\subsubsection{CEO Ownership: Lilienfeld-Toal and Ruenzi (2014)}

Lilienfeld-Toal and Ruenzi (2014), provide the cleanest portfolio based test of the individual level alignment hypothesis. They sort US firms annually into portfolios on lagged CEO ownership. They document a monotonic, statistically significant alpha that increases with ownership concentration. The high minus low ownership portfolio earns approximately 4-10\% annualised alpha under standard factor models across the board. Their preferred interpretation is that markets systematically under react to the long run benefits of managerial ownership at the individual decision maker level. This generates predictable abnormal returns to a portfolio strategy that buys high CEO ownership firms.

The Lilienfeld-Toal and Ruenzi (2014) study is the closest methodological precedent for this thesis. The portfolio sort architecture, the factor model adjustment, and the interpretive framing of alpha as evidence on the alignment hypothesis are all directly transferable. It is not, however, a perfect content equivalent. First, their identifying variable is CEO ownership. This thesis uses aggregate insider ownership, which directly measures the total stake of the firm's governance block. Second, their universe is restricted to ExecuComp-covered US firms. This excludes the European and Asia-Pacific firms that constitute a part of Wood's (2025) owner universe and which exhibit qualitatively different ownership institutions. Continental Europe in particular has more long tenured non founder insiders, and Asia-Pacific has concentrated family and associate ownership. 

These two differences are substantive rather than incremental. Aggregate insider ownership and CEO ownership measure different governance phenomena. A US only sample and a global sample cover different institutional environments. The pre 2010 and post 2015 periods reflect markedly different macroeconomic and capital market conditions. This thesis can therefore be positioned as adopting Lilienfeld-Toal and Ruenzi's methodology while studying a different ownership construct in a different geographic and time sample. Their core finding is a monotonic factor adjusted ownership alpha relationship. They generalize from CEO ownership to aggregate insider ownership in a global sample, which is an open empirical question that this thesis addresses directly.


\subsubsection{Family Ownership: Sraer and Thesmar (2007), Andres (2008), Lins et al. (2013)}

A third group of return relevant studies focuses on family ownership. Family ownership overlaps with but is conceptually different from CEO ownership. Sraer and Thesmar (2007) examine roughly 1000 French listed firms over 1994-2000. They document that family controlled firms generate superior accounting and stock return performance, with the effect concentrated in their managed and founder managed subgroups. Andres (2008) reports similar patterns in 275 German firms, with positive performance effects of founding family ownership conditional on active family involvement. Both studies use family firm indicator variables rather than continuous ownership measures. Both employ relatively limited factor model adjustment compared to Fahlenbrach (2009) or Lilienfeld-Toal and Ruenzi (2014).

Lins, Volpin, and Wagner (2013) provide an important counterweight. They are the only widely cited global study of ownership and returns. They examine 8854 firms across 35 countries during the 2008-2009 financial crisis. They report that family controlled firms underperformed non family firms by approximately 9\% during the crisis window. They attribute the underperformance to private benefit incentives when families face liquidity pressure. The Lins et al. (2013) result is important because it shows that under certain conditions, high insider or family ownership can predict negative rather than positive abnormal returns. However, the study is limited to a 24 month event window during a global crisis. It explicitly excludes US firms and uses a family control indicator rather than aggregate insider ownership. It is therefore informative about boundary conditions but is not a content precedent for this thesis.


\subsubsection{Aggregate Insider Ownership and Return Performance}

Bringing these interpretations together exposes a underdeveloped part in the literature that is not acknowledged. Of the more than one hundred studies surveyed by Wood (2025), only a handful examine realized stock returns rather than valuation ratios. Of these, the literature does not provide clear evidence on aggregate insider ownership as the sorting variable in a portfolio based factor model framework. Fahlenbrach (2009) studies founder CEO status. Lilienfeld-Toal and Ruenzi (2014) study CEO ownership. Sraer and Thesmar (2007) and Andres (2008) study family ownership in single European countries with limited factor adjustment. Lins et al. (2013) studies family control globally but only over a 24 month crisis window. The most comprehensive study of aggregate insider ownership is Wood (2025) and it is descriptive, with no factor model adjustment.

This research matters because aggregate insider ownership is the construct most directly aligned with the governance interpretation of the alignment versus entrenchment debate. When Morck, Shleifer, and Vishny (1988), Fama and Jensen (1983), and the broader theoretical literature discuss the conditions under which managerial alignment dominates entrenchment, the relevant decision unit is typically the firm's executive and board rather than any single individual. Studies that use CEO ownership alone capture only one member of that coalition. Studies that use family ownership alone capture only the subset of coalitions defined by family relationships. Studies that use founder CEO status capture only the subset in which the founder is still active. Aggregate insider ownership captures the full coalition and is therefore the most theoretically appropriate variable for testing collective alignment hypotheses.

This thesis therefore positions itself in the literature by adopting the portfolio sort and factor model methodology of Fahlenbrach (2009) and Lilienfeld-Toal and Ruenzi (2014). This is the established approach for risk adjusted return tests on ownership sorted portfolios. It applies this methodology to the universe assembled by Wood (2025), which is a comprehensive global sample of high insider ownership firms. It uses aggregate insider ownership rather than CEO, founder, or family ownership as the sorting variable, which addresses the part that the existing return literature has not directly examined.


\subsection{Cross Country Institutional Evidence and the Blockholder Mechanism}

Beyond the family ownership return studies discussed earlier in this paper, which create the basis of ownership literature, a broader cross country literature is analyzed. Maury and Pajuste (2005) study 136 Finnish firms. They find that the presence of multiple large blockholders dampen the entrenchment effects associated with single dominant owners. Barontini and Caprio (2006) extend the analysis to 675 firms across 11 continental European countries. They report a robust positive valuation effect of family control, measured through Tobin's Q rather than stock returns. These studies establish that the alignment entrenchment trade off documented in US data is qualitatively similar in continental European settings. The exact thresholds and effect sizes differ depending on the strength of minority shareholder protections and the presence of firms with multiple blockholders.

Evidence from concentrated ownership environments show the most important boundary condition for the alignment hypothesis. Claessens, Djankov, Fan, and Lang (2002) study eight East Asian economies. They document that the alignment benefits of large insider stakes hold only when voting rights are not significantly in excess of cashflow rights. Where dual class shares or pyramidal structures separate control from cashflow ownership, the relationship between insider control and firm value turns negative. Bennedsen and Nielsen (2010) replicate this finding in a sample of more than 4000 European firms. Wood (2025) explicitly name both papers in defending that he excluded dual class structures from his owner classification. He reports that, dual class firms substantially underperform single class high ownership firms. The implication for this thesis is that the Wood universe, which excludes dual class structures, represents the cleaner alignment case. Here insider voting and cashflow rights are aligned and the entrenchment via control rights mechanism is largely absent.

Another part that is examined in the literature is the broader role of large blockholders in corporate governance. It provides the theoretical mechanism in which aggregate insider ownership would be expected to affect long horizon performance. Edmans (2014), distinguishes between blockholder governance through "voice" (active intervention) and "exit" (informed trading). He argues that the empirical evidence on blockholder effects is highly sensitive to how blockholders are identified and what governance mechanisms are in operation. Insider blockholders are a special case in which the blockholders are simultaneously the firm's executives. This eliminates the voice versus exit distinction and concentrates both monitoring and decision making in the same coalition. Holderness, Kroszner, and Sheehan (1999) provide useful historical context. They document that managerial stock ownership in the US rose substantially through the twentieth century before reversing after 2006. Wood (2025) extends this trend and uses it to motivate the increasing distinctiveness of high ownership firms in contemporary markets.

Taken together, the cross country and blockholder literatures establish two propositions relevant to this thesis. First, the alignment versus entrenchment trade off appears qualitatively robust across institutional environments, conditional on the absence of voting rights versus cashflow rights divergence. Second, the existing cross country evidence is overwhelmingly value based rather than return based. The global return evidence outside the family ownership crisis study by Lins et al. (2013) is essentially limited to the descriptive analysis in Wood (2025). This thesis directly addresses this by applying a uniform asset pricing methodology to a genuinely global ownership sorted portfolio that excludes dual class structures.


\subsection{Asset Pricing Methodology: From CAPM to FF5}

The methodological backbone of this thesis is the Fama and French (2015) five factor model. This model extends the original Fama and French (1993) three factor specification by adding profitability (RMW) and investment (CMA) factors. The original three factor model augments the market excess return with size (SMB) and value (HML) factors. It was shown to capture the cross section of average US stock returns substantially better than the single factor CAPM. The five factor extension was motivated by valuation theoretic arguments and by empirical evidence that profitability and investment patterns produce cross sectional return spreads that are not absorbed by the original three factors. Carhart (1997) introduced a fourth momentum factor (UMD) that explains the return predictability associated with past 12 month returns. Many empirical studies augment the FF5 model with UMD as a robustness check.

In the context of ownership sorted portfolios, the FF5 framework is important because high ownership firms are systematically smaller and more value tilted than the broad market. Any apparent return advantage of high ownership portfolios under a single factor specification could therefore be entirely explained by SMB and HML. Apart from that the investment behaviour documented for high ownership firms by Fahlenbrach (2009) and Wood (2025) contains higher capital expenditure, lower payout, more reinvestment, which is exactly the type of behaviour that loads negatively on the CMA (conservative minus aggressive investment) factor. Failing to control for CMA would mechanically generate a positive alpha for high ownership portfolios that reflects their conservative investment rather than any genuine risk adjusted outperformance. This thesis therefore uses the global FF5 factors as the baseline specification, augmented with UMD as a robustness check, in order to ensure that any documented alpha is not absorbed by standard risk factor exposures.

A methodological refinement that has become standard in this literature is the use of Newey and West (1987) heteroskedasticity and autocorrelation consistent standard errors when estimating monthly portfolio regressions. Newey West correction is important because portfolio returns for portfolios with concentrated industry or country exposures can exhibit serial correlation that biases test statistics towards rejection of the null. This thesis follows Fahlenbrach (2009) and Lilienfeld-Toal and Ruenzi (2014) in applying Newey West corrections with a lag length appropriate for monthly data.


\subsection{Positioning Within the Literature: Wood (2025)}

Wood's (2025) study, is the guidance for this thesis. Wood assembles a global sample of 1297 firms with at least 20\% insider ownership and inflation adjusted market capitalisation of at least one billion US dollars. He tracks them over the period 1994-2023 for returns and 2013-2022 for fundamentals. He documents that this owner group generated annualized total returns of 10.9\% per year against 6.3\% per year for the S\&P Global 1200 benchmark. Beyond raw returns, Wood documents a series of behavioural differences between owner firms and the benchmark, including higher revenue growth, higher fixed asset investment, lower dividend and buyback payout, lower incidence of earnings guidance and higher returns on invested capital. He interprets these patterns as evidence that high insider ownership is associated with a long horizon stewardship orientation rather than short term profit maximization. His Figure 26 shows a monotonically positive relationship between inside ownership quintiles and annualized returns over five-, ten-, fifteen-, and twenty year horizons.

Wood's analysis is explicitly descriptive rather than econometric. He acknowledges that within his owner sample, the simple correlation between ownership level and stock performance is only 0.07. His attempts to characterise the ownership return relationship at finer ownership intervals yielded inconclusive results. He further notes that he conducts no risk adjustment, no factor model decomposition, and no formal test of statistical significance. These are not per se criticisms, but Wood frames his contribution as a descriptive analyses. That is precisely the limitation that this thesis is designed to address.

The underexplored part from this review can therefore be stated compactly. The cross sectional value literature is large, US dominated, and methodologically contested over the endogeneity question. The direct return literature is small, US dominated, and the literature is limited on aggregate insider ownership as a sorting variable. Fahlenbrach (2009) studies founder CEO status. Lilienfeld-Toal and Ruenzi (2014) study CEO ownership. Sraer and Thesmar (2007) and Andres (2008) study family ownership in single European countries. The only global return study (Lins et al., 2013) studies family control over a 24 month crisis window and excludes US firms. The most comprehensive global analysis of aggregate insider ownership (Wood, 2025) does not risk adjust returns, does not employ factor models, and does not test the statistical significance of the documented return differential. As far as the author of this thesis is aware there is no existing peer reviewed study which combines, aggregate insider ownership as the sorting variable, a global universe spanning North America, Europe, and Asia-Pacific, a long horizon extending through the post Covid period, and in combination with a Fama-French five factor adjustment on ownership sorted portfolios. This thesis contributes by sorting Wood's universe into aggregate insider ownership quintiles and estimating monthly FF5 alphas over the period January 2015 to January 2026. Including both equal weighted and value weighted portfolio constructions and with a momentum augmented six factor specification as a robustness check.


\subsection{Hypotheses}

The literature reviewed above generates predictions about the relationship between aggregate insider ownership and risk-adjusted stock returns. The convergence of interest hypothesis (Jensen and Meckling, 1976) predicts that higher inside ownership improves firm performance through reduced agency costs. The entrenchment hypothesis (Fama and Jensen, 1983) reverses this prediction at very high ownership levels. Stewardship theory (Donaldson and Davis, 1991) predicts long run outperformance through reinvestment rather than profit extraction. The descriptive evidence in Wood (2025) appears most consistent with alignment and stewardship, but his analysis cannot rule out that the documented return differential reflects standard risk-factor exposures. Leaving the question of genuine abnormal performance unanswered.

This thesis tests these predictions by estimating Fama-French five-factor alphas on portfolios sorted by aggregate insider ownership. The hypotheses progress from the basic question of whether the high-ownership universe earns positive alpha, to monotonicity, to specification and weighting robustness, and finally to two geographic hypotheses that exploit the global nature of the Wood (2025) sample. The geographic dimension matters because within a uniform methodology: Lilienfeld-Toal and Ruenzi (2014) is US only, Sraer and Thesmar (2007) and Andres (2008) are European, and Lins et al. (2013) covers only a 24-month crisis window. Five hypotheses follow.

H1 : Equal- and value weighted portfolios of firms with insider ownership exceeding 20\%, from the global universe assembled by Wood (2025), generate positive and statistically significant FF5 alphas over the period January 2015 to January 2026. 

H2 : FF5 alphas are monotonically increasing in ownership concentration when firms are sorted into ownership quintiles. This is consistent with the convergence of interest hypothesis and with the S\&P 1200 quintile pattern documented in figure 26 of Wood (2025). The entrenchment prediction of Fama and Jensen (1983), which would imply non monotonicity or negative alpha at the highest quintile, is given as a false alternative.

H3 : The ownership alpha relationship is robust to the inclusion of a momentum factor in a six factor specification and value- as well as equal-weighted portfolio construction. The momentum control addresses potential UMD loadings (Carhart, 1997). The dual weighting scheme separates results driven by small firms within each quintile from those an investor could realistically capture.

H4 : The positive ownership alpha relationship holds in each of four regional subsamples. North America, Europe, Asia-Pacific, and Emerging Markets ex-Asia (covering Latin America, Africa, and the Middle East). Estimated separately with regional FF5 factors. This follows from the cross-country evidence in Sraer and Thesmar (2007), Andres (2008), Maury and Pajuste (2005), Barontini and Caprio (2006), and Wood (2025). The Emerging Markets ex-Asia subsample is expected to have a smaller firm count than the three larger regions, so the strongest inference for H4 will come from North America, Europe, and Asia-Pacific, with the Emerging Markets ex-Asia result reported as a complementary test.

H5 : The North American ownership alpha exceeds the European ownership alpha. The directional prediction follows from Wood (2025), whose figure 3 documents that US insider ownership has declined substantially since 2006. As high insider ownership firms become scarcer in US capital markets, they become more differentiated from the benchmark and may command a larger premium. Confirmation suggests that scarcity in US capital markets amplifies the alignment premium. Rejection suggests the ownership effect is broadly invariant across the two markets.

H1 and H2 address the central question of whether aggregate insider ownership predicts risk adjusted returns. H3 ensures the result is not a result of model or weighting choices. H4 and H5 exploit the global Wood (2025) sample to test whether the ownership return relationship is invariant across regional markets or amplified in specific market structures.
